\documentclass{situationist-preprint}

\title{The Inseparability of Thick Concepts}
\msid{MTE-15-a}
\msfamily{MTE-15}
\msdiscipline{Metaethics}
\mstype{Principal manuscript}
\mskeywords{thick concepts; disentangling argument; shapelessness; separabilism; evaluative concepts; lewdness; expressivism}

\begin{document}
\maketitle

\begin{abstract}
Thick concepts, such as \emph{lewd}, \emph{cruel}, \emph{courageous} and \emph{tactful}, seem to describe and evaluate at once. Separabilists hold that each such concept factors into a descriptive concept, which fixes its extension, and an evaluation, typically thin, which is attached to that extension: \enquote{lewd} would mean \enquote{sexual, and bad for that reason}. We argue that thick concepts cannot be factored. We state the factoring thesis precisely and derive from it two consequences: that a competent user who does not share the evaluation could project the concept to new cases, and that the evaluation is the same thin evaluation in every application. Both are false of \emph{lewd}. The extension of \enquote{lewd} has no shape at the descriptive level, so that only someone who grasps the evaluative point can follow it, and the evaluation it carries is not a thin \enquote{bad} but a specific response that itself shapes the extension. We consider the main separabilist replies, including the determinable strategy of Elstein and Hurka and the pragmatic account of Väyrynen, and we state how the result bears on expressivism (MTE-03) and on Moral Twin Earth (MTE-13).
\end{abstract}

\section{The problem}

\citet{Williams1985} called concepts such as \emph{treachery}, \emph{brutality} and \emph{courage} thick: their application is guided by the world, so that there is agreement in their use among those who share them, and yet they are action-guiding, so that to apply them is characteristically to have a reason for or against something. Thin concepts, such as \emph{good}, \emph{wrong} and \emph{ought}, are action-guiding but have little world-guided content. The question is whether a thick concept is a combination of the two kinds of content.

\emph{Separabilists} answer yes. \citet{Hare1952,Hare1963} held that evaluative words have a descriptive meaning, which gives the criteria for their application, and an evaluative meaning, which is the commendation or condemnation that they express, and that the two may come apart: one may use \enquote{rude} with its descriptive meaning only, in an inverted-commas sense. \citet{Blackburn1992,Blackburn1998} defended a version on which a thick term is a descriptive predicate with an attitude attached, so that the thick concept can be disentangled into a description and an attitude towards things so described. \citet{Gibbard1992} offered a related account in his exchange with Blackburn. \citet{ElsteinHurka2009} gave the most developed recent version: a thick concept combines a thin evaluation with a descriptive content that may be specified only as a determinable, whose determinates are fixed by the evaluation.

\emph{Inseparabilists} answer no. \citet{McDowell1981}, drawing on Wittgenstein's remarks on rule-following, argued that there need be no descriptive concept that captures the extension of a thick concept, since the extension may be shapeless from any standpoint that does not share the evaluative interest, and that a person who could not see the point of the concept could not go on correctly with it. \citet{Foot1958} had earlier argued that some evaluative terms, such as \enquote{rude}, cannot be applied without standards of application that make the evaluation non-optional. \citet{Dancy1995}, \citet{Putnam2002} and \citet{Roberts2013} have defended versions of inseparability.

We argue for inseparability, on the example that has become standard: \enquote{lewd} is not \enquote{sexual} plus \enquote{bad} \citep{Blackburn1992,Vayrynen2013}. The argument is general, but the example makes each step concrete.

\section{Framework}

\begin{definition}[Descriptive concept]\label{d:desc}
A concept is \emph{descriptive} if a thinker can master it, that is, apply it correctly to new cases, without sharing or understanding any evaluative attitude.
\end{definition}

\begin{definition}[Factoring]\label{d:factor}
A thick concept $T$ \emph{factors} if there is a descriptive concept $D$ and an evaluation $E$ such that (a) necessarily, something falls under $T$ if and only if it falls under $D$ and $E$ applies to it in virtue of its falling under $D$; and (b) mastery of $T$ consists in mastery of $D$ together with the disposition to make the evaluation $E$ of things that fall under $D$.
\end{definition}

The factoring thesis for $T$ is the claim that $T$ factors. In the simple form Hare and Blackburn defended, $E$ is a thin evaluation such as \emph{bad}, and the factoring of \emph{lewd} is \enquote{sexual, and bad because sexual}. In the more sophisticated form of \citet{ElsteinHurka2009}, $D$ may be a more specific descriptive concept, such as \enquote{sexual display exceeding certain limits}, and $E$ thin. The definition covers both. Clause (b) is what distinguishes factoring from mere necessary coextension: the separabilist claims not only that the extension of $T$ is the intersection of a descriptive extension and an evaluative one, but that competence with $T$ is decomposable accordingly.

\begin{principle}[Projection]\label{p:proj}
If a concept $D$ is descriptive, a thinker who has mastered $D$ can determine, for any new case, whether it falls under $D$, without sharing any evaluative attitude.
\end{principle}

\begin{principle}[Shapelessness]\label{p:shape}
There is no descriptive concept whose extension coincides, across the cases that competent users of \enquote{lewd} would classify, with the extension of \emph{lewd}, or with the extension of any concept $D$ such that \emph{lewd} is $D$ plus an evaluation applying in virtue of $D$.
\end{principle}

Principle~\ref{p:proj} follows from Definition~\ref{d:desc}. Principle~\ref{p:shape} is the substantive premise, and we defend it in Section~\ref{s:cases}.

\section{The argument}

\begin{lemma}\label{l:outsider}
If \emph{lewd} factors, a thinker who does not share the evaluative attitude that \emph{lewd} carries can sort new cases into those that fall under the descriptive component and those that do not.
\end{lemma}

\begin{proof}
If \emph{lewd} factors, there is a descriptive $D$ as in Definition~\ref{d:factor}. By Principle~\ref{p:proj}, a thinker who has mastered $D$ can sort new cases by $D$ without sharing any evaluative attitude. Mastery of $D$ is available to the outsider, since by Definition~\ref{d:desc} it requires no evaluative attitude.
\end{proof}

\begin{theorem}\label{t:main}
\emph{Lewd} does not factor.
\end{theorem}

\begin{proof}
Suppose it does, with descriptive component $D$ and evaluation $E$. By clause (a) of Definition~\ref{d:factor}, the extension of \emph{lewd} is the part of the extension of $D$ to which $E$ applies in virtue of $D$. Two cases arise. If $E$ applies to everything that falls under $D$ in virtue of its falling under $D$, then the extension of \emph{lewd} is the extension of $D$, and this contradicts Principle~\ref{p:shape} directly. If $E$ applies to some things under $D$ and not to others, then whether it applies in virtue of $D$ cannot depend on $D$ alone, since what falls under $D$ in virtue of a single feature receives the evaluation uniformly; so it depends on features beyond $D$, and the concept $D'$ of falling under $D$ and having those features is either descriptive, in which case it is a descriptive concept whose extension coincides with that of \emph{lewd} and Principle~\ref{p:shape} is again contradicted, or not descriptive, in which case the outsider of Lemma~\ref{l:outsider} cannot determine which cases under $D$ are lewd, and clause (b) of Definition~\ref{d:factor} fails, since mastery of $D$ together with the evaluation $E$ does not settle the extension. In every case the supposition fails.
\end{proof}

The proof has the form of the disentangling argument that \citet{McDowell1981} directed against non-cognitivism, but it does not depend on the further claims that McDowell drew from it about moral perception. It needs only Principle~\ref{p:shape} and the definition of factoring. The second case of the proof is the one that defeats the sophisticated separabilist: if the description must be completed by features that only the evaluation can pick out, the description has not been separated from the evaluation.

\section{Cases}\label{s:cases}

Principle~\ref{p:shape} is tested against the descriptive candidates, each paired with a case that defeats it.

\paragraph{Sexual.} A medical diagram of the reproductive system is sexual in content, and nobody thinks it lewd. A married couple's private intimacy is sexual, and is not lewd. A sexual assault is sexual and bad, and \enquote{lewd} is the wrong word for it: it understates the wrong and misdescribes its character. So \enquote{sexual and bad} is too broad on the bad side as well as the descriptive side.

\paragraph{Sexually explicit.} A frank but tender love poem may be explicit without being lewd, and a glance or innuendo may be lewd without being explicit. Explicitness and lewdness come apart in both directions.

\paragraph{Sexual display in public.} The candidate improves on the first two, since lewdness has to do with display. But a public kiss at a wedding is a sexual display in public and is not lewd, while a remark made in private to a colleague may be lewd. The candidate also leaves out what lewdness consists in: not publicity as such, but a display that obtrudes sexuality where it is unwelcome or out of place, in a way that degrades the occasion or the person addressed.

\paragraph{Sexual display that exceeds the limits of decency.} This candidate, close to Blackburn's own gloss, at last tracks the extension. But \enquote{decency} is itself thick. The candidate is not descriptive in the sense of Definition~\ref{d:desc}, since one cannot master the limits of decency without understanding the evaluative interest in sexual modesty that sets them. Replacing \enquote{decency} by a list of conventions, \enquote{sexual display that violates the conventions of the community}, makes the candidate descriptive but wrong: conventions may be prudish, so that a display that violates them is not lewd, and a community may tolerate lewdness, so that a lewd display violates no convention.

The pattern is the one Principle~\ref{p:shape} states. Each descriptive candidate is either too broad or too narrow, and the candidates that get the extension right do so by importing an evaluative concept. The shape of the extension is visible only from within the evaluative interest that gives the concept its point.

\section{The evaluation is not thin}

The separabilist may grant that the description cannot be specified independently and still hold that the evaluation is a separable thin \emph{bad}. That is also false, for two reasons.

First, the evaluation is specific. To call a remark lewd is not merely to condemn it but to condemn it in a particular way, as obtrusive and degrading in its sexuality; to call it cruel is to condemn it in another way. Two remarks may be equally bad, one lewd and the other cruel, and the difference between the concepts is not a difference in what is condemned on the descriptive side together with a common \emph{bad}, since we have seen that there is no descriptive side to be had. The difference is in the evaluation itself.

Second, the evaluation is defeasible in a way thin \emph{bad} is not. A speaker may say \enquote{it was lewd, and wonderful}, of a music-hall song. \citet{Blackburn1992} took such cases to show that the evaluation is separable, since here the description applies and the condemnation does not. \citet{Vayrynen2013} took them to show that the evaluation is not part of the meaning at all but a pragmatic implication. We read the case differently. The speaker who calls the song lewd and wonderful applies the concept with its evaluative point and celebrates the transgression it identifies: the song is wonderful partly because it is lewd, because it flouts the standards of modesty that \enquote{lewd} invokes. The evaluation is present, and it is overridden by a further evaluation of transgression on that occasion. Thin \emph{bad} does not behave this way, since nothing is wonderful partly because it is bad. The evaluation that lewdness carries is not thin.

\section{Objections and replies}

\begin{objection}[The determinable strategy]
\citet{ElsteinHurka2009} grant that the descriptive content of a thick concept may be specifiable only as a determinable: \enquote{sexual display with properties that make it bad}. The determinable is descriptive, the evaluation is thin, and the concept factors.
\end{objection}

\begin{reply}
A determinable of the form \enquote{having some descriptive properties that make it bad} is descriptive only in name. Clause (b) of Definition~\ref{d:factor} requires that mastery of the descriptive component, together with the evaluation, settle the extension; but mastery of the determinable does not tell the outsider which determinates fall under it, since that is fixed by the evaluation. The determinable strategy is the second case of the proof of Theorem~\ref{t:main}: the description is completed by features that only the evaluation picks out. Elstein and Hurka's analysis may correctly state a necessary equivalence, but it is not a factoring.
\end{reply}

\begin{objection}[Pragmatic evaluation]
The evaluation is not part of the meaning of \enquote{lewd} but a presupposition or implicature that attaches to its use in typical contexts \citep{Vayrynen2013}. If so, the concept is purely descriptive, and the question of factoring does not arise.
\end{objection}

\begin{reply}
If \enquote{lewd} were descriptive, Principle~\ref{p:shape} would be false, since \emph{lewd} itself would be a descriptive concept with the extension of \emph{lewd}. The pragmatic account must therefore hold either that some purely descriptive concept captures the extension, which the cases of Section~\ref{s:cases} tell against, or that an outsider could master \enquote{lewd} without understanding the evaluation it typically conveys. The second option is implausible: an outsider who had learned to apply \enquote{lewd} to the cases on which competent users apply it, and who did not understand why they cared about the cases so classified, would not be able to classify new ones. Whether the evaluation is encoded semantically or pragmatically, it is required for mastery of the extension, and that is what inseparability claims.
\end{reply}

\begin{objection}[Anthropologists]
Anthropologists routinely learn the thick concepts of cultures whose values they do not share, and apply them correctly. So the outsider can master thick concepts.
\end{objection}

\begin{reply}
The anthropologist succeeds by coming to understand the evaluative interest of the culture, imaginatively entering its point of view, without endorsing it. That is consistent with inseparability, which says that mastery requires grasping the evaluation, not sharing it. What the anthropologist cannot do is master the concept by learning a descriptive classification and then attaching an attitude; her success depends on her understanding of what the people she studies care about. Definition~\ref{d:desc} counts a concept as descriptive only if it can be mastered without understanding any evaluative attitude, and the anthropologist's case shows that thick concepts are not descriptive in this sense.
\end{reply}

\section{Relations to other results}

\paragraph{Expressivism.} Separabilism has been the expressivist's account of thick concepts, since it allows the descriptive component to be represented by belief and the evaluation by attitude \citep{Blackburn1992}. MTE-03 argues that expressivism can solve the Frege--Geach problem by taking every sentence to express a state of being for a set of fact-plan worlds, and shows that such contents need not be intersections of a descriptive and a practical content. An expressivist who accepts that semantics can therefore accept the present conclusion: she may assign to \enquote{$x$ is lewd} a content that is not the intersection of a descriptive and a practical one, as \citet{Gibbard1992} in effect proposed. Inseparability is consistent with expressivism of the MTE-03 kind. What it rules out is the traditional expressivist explanation of thick terms as descriptive predicates with attitudes attached.

\paragraph{Realism and Moral Twin Earth.} Inseparability is also consistent with the realism of MTE-01, on which moral facts are natural facts about idealized reactions. The property that \emph{lewd} picks out may be a natural property, if it is a property of eliciting a specific response under idealized conditions; what Theorem~\ref{t:main} denies is that it is the conjunction of an evaluatively neutral descriptive property and a thin evaluation. The specific response is part of what lewdness is, not a thin verdict attached to an independently given description. MTE-13 argues that moral terms refer to the same properties on Twin Earth, so that apparent divergence in use reflects error rather than difference of meaning. Inseparability supports the conclusion for thick terms: since the extension of \emph{lewd} is fixed by the evaluative interest it serves and not by a descriptive classification, a community that applies its counterpart of \enquote{lewd} to a different range of cases, while sharing the interest, is misapplying the same concept.

\paragraph{Hume's law.} MTE-10 argues that normative conclusions follow from purely descriptive premises only with a normative bridge premise. Since \emph{lewd} is not descriptive, an argument that concludes that a remark is lewd needs a premise that is not purely descriptive, whether a thin principle or another thick judgment. This is the thick counterpart of Hume's law, and it explains why disputes about whether something is lewd cannot be settled by agreement on the descriptive facts alone.

\section{Conclusion}

\enquote{Lewd} is not \enquote{sexual} plus \enquote{bad}. Every descriptive candidate for the descriptive component either misclassifies cases or imports an evaluative concept, so that the extension of \emph{lewd} has no shape at the descriptive level, and the evaluation it carries is specific and shapes the extension rather than being a thin verdict added to it. Thick concepts cannot be factored into description and evaluation.

\begin{thebibliography}{99}\small
\bibitem[Blackburn(1992)]{Blackburn1992} Blackburn, S. (1992). Morality and thick concepts: Through thick and thin. \emph{Proceedings of the Aristotelian Society, Supplementary Volume}, 66, 285--299.
\bibitem[Blackburn(1998)]{Blackburn1998} Blackburn, S. (1998). \emph{Ruling Passions: A Theory of Practical Reasoning}. Clarendon Press.
\bibitem[Dancy(1995)]{Dancy1995} Dancy, J. (1995). In defense of thick concepts. \emph{Midwest Studies in Philosophy}, 20, 263--279.
\bibitem[Elstein and Hurka(2009)]{ElsteinHurka2009} Elstein, D.~Y. and Hurka, T. (2009). From thick to thin: Two moral reduction plans. \emph{Canadian Journal of Philosophy}, 39(4), 515--535.
\bibitem[Foot(1958)]{Foot1958} Foot, P. (1958). Moral arguments. \emph{Mind}, 67(268), 502--513.
\bibitem[Gibbard(1992)]{Gibbard1992} Gibbard, A. (1992). Morality and thick concepts. \emph{Proceedings of the Aristotelian Society, Supplementary Volume}, 66, 267--283.
\bibitem[Hare(1952)]{Hare1952} Hare, R.~M. (1952). \emph{The Language of Morals}. Clarendon Press.
\bibitem[Hare(1963)]{Hare1963} Hare, R.~M. (1963). \emph{Freedom and Reason}. Clarendon Press.
\bibitem[McDowell(1981)]{McDowell1981} McDowell, J. (1981). Non-cognitivism and rule-following. In S. Holtzman and C. Leich (eds.), \emph{Wittgenstein: To Follow a Rule}, 141--162. Routledge \& Kegan Paul.
\bibitem[Putnam(2002)]{Putnam2002} Putnam, H. (2002). \emph{The Collapse of the Fact/Value Dichotomy and Other Essays}. Harvard University Press.
\bibitem[Roberts(2013)]{Roberts2013} Roberts, D. (2013). Thick concepts. \emph{Philosophy Compass}, 8(8), 677--688.
\bibitem[V\"ayrynen(2013)]{Vayrynen2013} V\"ayrynen, P. (2013). \emph{The Lewd, the Rude and the Nasty: A Study of Thick Concepts in Ethics}. Oxford University Press.
\bibitem[Williams(1985)]{Williams1985} Williams, B. (1985). \emph{Ethics and the Limits of Philosophy}. Fontana.
\end{thebibliography}

\end{document}
